\documentclass[10pt,a4paper]{amsart}
\usepackage[margin=19mm]{geometry}
\usepackage{amsmath,amssymb,mathtools}
\usepackage[scr=rsfs,cal=euler]{mathalfa}
\usepackage{xfrac}
\usepackage[unicode,hidelinks,bookmarks=false]{hyperref}
\newcommand{\ie}{i.e.\ }
\newcommand{\cf}{cf.\ }
\newcommand{\resp}{respectively\ }
\newcommand{\Subsection}{\subsection}
\providecommand{\nobreakdash}{\nobreak\hbox{-}}
\setlength{\emergencystretch}{2em}
\multlinegap0pt
\usepackage{siunitx}
\usepackage{siunitx}
\usepackage{siunitx}
\usepackage{siunitx}
\usepackage{siunitx}
\usepackage{tikz}
\newtheorem{lemma}{Lemma}
\usepackage{cleveref}
\crefname{lemma}{Lemma}{Lemmas}
\newcommand{\SU}{{\mathcal U}}
\title{Geometric Conditions for Lossless Convexification in Linear Optimal Control with Discrete-Valued Inputs: Real-Time Implementation for Spacecraft Rendezvous}
\author{Felipe Arenas-Uribe, Hasan A. Poonawala, Jesse B. Hoagg}
\date{Source: arXiv:2511.07711v3 (CC BY 4.0)}
\begin{document}
\maketitle

\noindent\emph{Source and licence.} Geometric Conditions for Lossless Convexification in Linear Optimal Control with Discrete-Valued Inputs: Real-Time Implementation for Spacecraft Rendezvous. Version arXiv:2511.07711v3 (CC BY 4.0), distributed by arXiv under the Creative Commons Attribution 4.0 International licence, http://creativecommons.org/licenses/by/4.0/, as recorded in the arXiv licence metadata for this exact version and shown on the abstract page https://arxiv.org/abs/2511.07711v3. Full licence notice: \texttt{licenses/arxiv-2511.07711-CC-BY-4.0.txt}.\par\smallskip

\noindent\textit{Editorial scope.} Counted unit: the article's lemma lemma:MICP2mayer (source0.tex lines 291--293). The article's complete original proof of it is delivered with it (source0.tex lines 295--325, one \texttt{proof} environment of 31 lines). No mathematics is written, repaired or re-derived: the statement and the proof are the article's own lines, in the article's own notation, and no retained line is substituted or re-flowed.  Retained uncounted as the source-local context the proof needs: lines 265--267; lines 283--285.  Nothing was omitted from any retained span, so the package carries no omission record.  No retained line contains a citation, so the wrapper carries no bibliography and no bibliography entry.  No retained line contains a figure environment, an image-inclusion command or drawing code, so no image is delivered and the package declares no asset dependency.  Editorial formatting only, in addition: an AMS \texttt{amsart} preamble that restates the article's own declarations and macros used by the retained lines, namely the environments \texttt{lemma} on separate counters, together with the standard packages those lines need.  The retained environments print in this document as Lemma 1 in order of appearance, whereas the article prints the same items with the numbering of its own full document.  The complete native source of the article as distributed in the arXiv e-print for this exact version is bundled unchanged as \texttt{sources/188831/source0.tex}.
\begin{equation}\label{eq:optimal_control_p1}
    u^*(t) = \arg \min_{u \in \SU} \big[ H_1(t,x,u,p) \big], \quad \forall t \in [0,t_f].
\end{equation}

\begin{equation}\label{eq:optimal_control_p2}
    \hat u^*(t) = \arg \min_{\hat u \in \SU_e} \big[ H_2(t,\hat x,\hat u,p, p_c) \big], \quad \forall t \in [0,t_f].
\end{equation}

\begin{lemma}\rm\label{lemma:MICP2mayer}
A control $u^*(t):[0,t_f]\to\mathbb{R}^m$ minimizes \cref{eq:optimal_control_p1} if and only if the pair $\big(u^*(t),\mu^*(t)\big)$, where $\mu^*(t)=\psi(u^*(t))$, minimizes \cref{eq:optimal_control_p2}. 
\end{lemma}

\begin{proof}
    Fix $t\in[0,t_f]$. By definition of the augmented input set $\mathcal{U}_e$, every admissible pair $(u,\mu)\in\mathcal{U}_e$ satisfies
    %
    \begin{equation*}
    \mu \ge \psi(u).
    \end{equation*}
    %
    Then, for all $(u,\mu) \in \SU_e$,
    %
    \begin{equation*}
    p(t)^\top\big(A(t)x(t)+B(t)u(t)\big) + \mu \geq p(t)^\top\big(A(t)x(t)+B(t)u(t)\big) + \psi(u).
    \end{equation*}
    %
    Equality holds if and only if $\mu=\psi(u)$. Which implies that for all $(u,\mu) \in \SU_e$,
    %
    \begin{equation*}
    H_2(t,\hat x,\hat u,p, p_c) \geq H_1(t,x,u,p).
    \end{equation*}
    
    It follows that
    %
    \begin{equation*}
    \min_{(u,\mu)\in\mathcal{U}_e} H_2(t,\hat x,\hat u,p,p_c)
    =
    \min_{u\in\mathcal{U}} H_1(t,x,u,p),
    \end{equation*}
    %
    and the minimizers satisfy $\mu^*=\psi(u^*)$.
    
    Therefore, $u^*(t)$ minimizes \Cref{eq:optimal_control_p1} if and only if the pair $\big(u^*(t),\mu^*(t)\big)$ with $\mu^*(t)=\psi(u^*(t))$ minimizes \Cref{eq:optimal_control_p2}.
\end{proof}

\end{document}
